\documentclass{beamer}
\usetheme{JuanLesPins}
\usecolortheme{seahorse}

\usepackage{fancyvrb}
\usepackage{graphicx}
\usepackage[style=authortitle-icomp]{biblatex}



%gets rid of bottom navigation bars
\setbeamertemplate{footline}[]{}

%gets rid of navigation symbols
\setbeamertemplate{navigation symbols}{}

\bibliography{thesis}

\title{Push-Pull Signal-Function Functional Reactive Programming}
\author{Edward Amsden}
\institute{Rochester Institute of Technology}

\begin{document}

\AtBeginSection[]
{
\begin{frame}
\frametitle{\insertsection}
\end{frame}
}

\begin{frame}[plain]
\titlepage
\includegraphics[width=9em]{RIT_GCCIS1}
\end{frame}


\section{Introduction}

\begin{frame}
\frametitle{Functional Reactive Programming (FRP)}
    \begin{itemize}
        \item {\em Functional} First-class and higher-order functions.
        \item {\em Reactive} Behavior changes in response to temporal inputs.
        \item Basic abstractions
        \begin{itemize}
            \item Signals - Functions of time\footnote{Often also called behaviors.}.
            \item Events - Sequences of temporally ordered and labeled discrete values.
        \end{itemize}
    \end{itemize}
\end{frame}

\begin{frame}
\frametitle{Signal-Function FRP}
    \begin{itemize}
        \item Signal functions are transformers of events and signals.
        \item Signal functions are first class in Signal-Function FRP.
        \item Signals and events are not first-class in Signal Function FRP.
        \item This approach is more composable than first-class signals and events (since input may be transformed)\footcite{Courtney2001-1}.
    \end{itemize}
\end{frame}

\begin{frame}
\frametitle{Push vs. Pull Evaluation}
    \begin{itemize}
        \item When to evaluate what?
        \item Pull evaluation: (``demand-driven'') Evaluate when output is needed.
        \item Push evaluation: (``data-driven'') Evaluate when input is available.
        \item Ideally, FRP uses both: pull for signals and push for events.
    \end{itemize}
\end{frame}

\section{Problem}

\begin{frame}[fragile]
\frametitle{Separating Events and Signals}
    \begin{itemize}
        \item[] 
\begin{Verbatim}
class (Category a) => Arrow a where
  arr :: (b -> c) -> a b c
  ...
\end{Verbatim}
        \item Traditional Signal-Function FRP encodes signal functions as a Haskell Arrow.
\begin{Verbatim}
data SF a b = ...

instance Arrow SF where
  ...
\end{Verbatim}
    \end{itemize}
\end{frame}

\begin{frame}[fragile]
\frametitle{Separating Events and Signals}
    \begin{itemize}
        \item Necessary to lift any function to a corresponding signal function, without input/output type annotation
        \item Events encoded as optional signals:
\begin{Verbatim}
type Event a = Maybe a
\end{Verbatim}
        \item Combined signals and events encoded as one signal with pair type:
    \end{itemize}
\end{frame}

\begin{frame}[fragile]
\frametitle{Separating Events and Signals}
    \begin{itemize}
        \item[]
\begin{Verbatim}
class (Category a) => Arrow a where
  ...
  first :: a b c -> a (b, d) (c, d)
  ...

capture :: SF (a, Event b) (Event a)

ghci> :t first . arr
first . arr
    :: (Arrow a) => (b -> c) -> a (b, d) (c, d)
\end{Verbatim}
    \end{itemize}
\end{frame}

\section{Approach}

\begin{frame}[fragile]
\frametitle{Type Signal Functions with Signal Vectors}

    \begin{itemize}
        \item Described by Sculthorpe\footcite{Sculthorpe2011}
        \item Describe separation of individual signals and events
        \item In Haskell (with {\tt -XEmptyDataDecls}):
\begin{Verbatim}
data SVEmpty
data SVSignal a
data SVEvent a
data SVAppend svl svr
type :^: svl svr = SVAppend svl svr
\end{Verbatim}
        \item Would like to use {\tt -XDataKinds} but it's not working well yet.
    \end{itemize}

\end{frame}

\begin{frame}[fragile]
\frametitle{Type Signal Functions with Signal Vectors}

    \begin{itemize}
        \item Two combinator examples:
        \item Lifting a pure function to transform a signal:
\begin{Verbatim}
pureSignal ::     (a -> b)
               -> SF (SVSignal a) (SVSignal b)
\end{Verbatim}
        \item Passing through input on the right:
\begin{Verbatim}
first ::    SF svIn svOut
         -> SF (svIn :^: svRight) (svOut :^: svRight)
\end{Verbatim}
        \item Composing these leads to:
\begin{Verbatim}
ghci> :t first . pureSignal
first . pureSignal
    :: (a -> b) -> SF (SVSignal a :^: sv)
                      (SVSignal b :^: sv)
\end{Verbatim}
    \end{itemize}
\end{frame}

\begin{frame}
\frametitle{Partial Representations of Signal Vectors}
    \begin{itemize}
        \item To evaluate signals and events differently we must represent them separately.
        \item Signal vectors enable this by distinguishing them in the types.
        \item We can construct several representations of a signal vector.
        \begin{itemize}
          \item Represent signal leaves, event leaves, or both.
          \item Represent one leaf, a subset of leaves, or all applicable leaves.
          \item Transform the type at the leaf, or don't.
        \end{itemize}
    \end{itemize}
\end{frame}

\section{Implementation}

\begin{frame}[fragile]
\frametitle{Signal Representation}
    \begin{itemize} 
        \item Represent an entire signal sample\\ (for initializing a signal function at time zero):
\begin{Verbatim}
data SVSample where
    SVSample      ::    a -> SVSample (SVSignal a)
    SVSampleEvent ::    SVSample (SVEvent a)
    SVSampleEmpty ::    SVSample SVEmpty
    SVSampleBoth  ::    SVSample svLeft 
                     -> SVSample svRight
                     -> SVSample (svLeft :^: svRight) 
\end{Verbatim}
    \end{itemize}
\end{frame}

\begin{frame}[fragile]
\frametitle{Signal Representation}
    \begin{itemize} 
        \item Represent a signal delta\\ (replacement values for a subset of signals):
\begin{Verbatim}
data SVDelta where
     SVDelta        ::   a -> SVDelta (SVSignal a)
     SVDeltaNothing ::   SVDelta sv
     SVDeltaBoth    ::   SVDelta svLeft 
                      -> SVDelta svRight
                      -> SVDelta (svLeft :^: svRight) 
\end{Verbatim}
    \end{itemize}
\end{frame}

\begin{frame}[fragile]
\frametitle{Event Representation}
    \begin{itemize} 
        \item Represent an event occurrence:
\begin{Verbatim}
data SVOccurrence where
     SVOccurrence ::   a -> SVOccurrence (SVEvent a)
     SVOccLeft    ::   SVOccurrence svLeft
                    -> SVOccurrence
                         (svLeft :^: svRight)
     SVOccRight   ::   SVOccurrence svRight
                    -> SVOccurrence
                         (svLeft :^: svRight)
\end{Verbatim}
    \end{itemize}
\end{frame}

\begin{frame}[fragile]
\frametitle{Signal Function Representation}
    \begin{itemize}
        \item Separate continuations for time advancement and event occurrences:
\begin{Verbatim}
data Initialized
data NonInitialized

data SF init svIn svOut where
  SF ::   (SVSample svIn
           -> (SVSample svOut,
               SF Initialized svIn svOut))
       -> SF NonInitialized svIn svOut
...
\end{Verbatim}
    \end{itemize}
\end{frame}

\begin{frame}[fragile]
\frametitle{Signal Function Representation}
    \begin{itemize}
    \item Separate continuations for time advancement and event occurrences.
\begin{Verbatim}
...
  SFInit ::   (Double 
               -> SVDelta svIn
               -> (SVDelta svOut,
                   [SVOccurrence svOut],
                   SF Initialized svIn svOut))
           -> (SVOccurrence svIn
               -> ([SVOccurrence svOut],
                   SF Initialized svIn svOut)
           -> SF Initialized svIn svOut
\end{Verbatim}
    \end{itemize}
\end{frame}

\begin{frame}
\frametitle{Evaluation}
    \begin{itemize}
        \item Yampa/AFRP: Supply SF and input/output actions to an evaluation loop ({\tt reactimate}).
        \item Here: initialize an evaluation state with:
        \begin{itemize}
            \item A signal function.
            \item Initial values for all input signals.
            \item Handlers for all outputs.
        \end{itemize}
        \item Then, the evaluation monad carries this state and provides the actions:
            \begin{itemize}
                \item {\tt push} Push an event (which will be immediately reacted to).
                \item {\tt update} Update the value of an input signal (with no immediate effect).
                \item {\tt step} Update the time and evaluate new values of signals.
            \end{itemize}
    \end{itemize}
\end{frame}

\section{Conclusion}

% Add conclusion

\begin{frame}
\frametitle{Further Work}
\begin{itemize}
  \item Dynamic collections: Dynamically switch between collections of signal functions.
  \item Semantics/correctness proof (especially for event merging).
  \item Time-independence optimization.
\end{itemize}
\end{frame}

\begin{frame}
\frametitle{Questions?}
\end{frame}

\section{In Detail}

\begin{frame}
\frametitle{Combinator Implementations: Serial Composition ({\tt >>>})}
    \begin{itemize}
        \item Initialization:
        \item Running:
    \end{itemize}
\end{frame}

\begin{frame}
\frametitle{Combinator Implementations: Parallel Composition ({\tt first})}
    \begin{itemize}
        \item Initialization:
        \item Running:
    \end{itemize}
\end{frame}

\begin{frame}
\frametitle{Combinator Implementations: Parallel Composition ({\tt second})}
    \begin{itemize}
        \item Initialization:
        \item Running:
    \end{itemize}
\end{frame}

\begin{frame}
\frametitle{Time-independence optimization}
\end{frame}

\begin{frame}
\frametitle{Event Merging}
\end{frame}

\end{document}
